\documentclass[11pt]{article}
\usepackage{icmlko}
\begin{document}

\title{고르개를 고쳤는데도 진다 --- 문제는 고르개가 아니라 안쪽 시기였다}
\author{월드모델 실험실}
\date{노트 499 · 2026년 8월 3일}
\maketitle

\begin{abstract}
노트 498 이 $1$ 차 고르개의 패인을 진단했다 --- \emph{고르개가 이웃 값을
못 본다}(도메인별 점수만 보고 고르는데, 공유 축 모형에서는 한 칸을 바꾸면
모든 칸이 바뀐다). 그래서 고쳤다. 후보 $21$ 짝을 하나씩만 가리고
\textbf{안쪽 판}(이웃 포함, 무게까지 반영)을 쟀다.

\textbf{고르개는 실제로 고쳐졌다.} 작은 도메인들의 안쪽 판 차가 제자리로
내려앉는다 --- 팝업 $+0.0003$ · 아이돌 $+0.0010$ · 게임 $+0.0000$ ·
시장팝업 $+0.0000$. $1$ 차에서 이들의 도메인 점수 차가 $+0.07 \cdot
+0.05 \cdot +0.009$ 로 제일 커 보였던 것과 정반대다. 뽑힌 표도 합리적이다
--- 애니 트렌드 · 모바일 $\{$위키, 태그, gen$\}$ · 팝업 media\_push,
\textbf{합친 안쪽 판 $+0.0083$}($3/3$).

\textbf{그런데 유보에서 진다 --- 판 $-0.0078$ · $0/12$.}
\end{abstract}

\begin{figure}[t]
\centering
\includegraphics[width=\textwidth]{figs/f1.pdf}
\caption{\textbf{가} --- 스물한 짝의 안쪽 판 차. \textbf{나} --- 범인
가르기. \textbf{다} --- 고르개 셋.}
\end{figure}

\section{범인은 모바일 셋}

노트 498 채택 상태에 하나씩 더해 봤다.

\begin{center}
\begin{tabular}{lrrr}
\hline
팔 & 판 차 & 부호 & 날짜 통제 차 \\
\hline
$+$팝업 media\_push & $+0.0004$ & $9/12$ & $+0.0009$ \\
$+$모바일 $\{$위키, 태그, gen$\}$ & $\mathbf{-0.0073}$ & $0/12$ & $-0.0098$ \\
둘 다 & $-0.0078$ & $0/12$ & $-0.0098$ \\
\hline
\end{tabular}
\end{center}

\textbf{더해진다}($-0.0069$ 대 실제 $-0.0078$). 안쪽에서는 합칠 때
줄어들었는데($+0.0094 \to +0.0083$) 유보에서는 안 줄어든다.

\section{크기가 답을 안다}

안쪽 판 차를 크기순으로 세우면 틈이 있다.

\begin{center}
\begin{tabular}{llr}
\hline
도메인 & 축 & 안쪽 판 차 \\
\hline
\textbf{애니} & \textbf{트렌드} & $\mathbf{+0.0042}$ \\
모바일 & 위키 & $+0.0019$ \\
모바일 & gen & $+0.0019$ \\
모바일 & 태그 & $+0.0011$ \\
팝업 & media\_push & $+0.0003$ \\
\hline
\end{tabular}
\end{center}

\textbf{$1$ 등이 $2$ 등의 두 배가 넘는다.} 그리고 유보에서 이기는 것은
$1$ 등 하나뿐이다. 문턱을 $(0.0019,\ 0.0042)$ 어디에 놓아도 \textbf{노트
498 이 채택한 바로 그 표}가 나온다 --- 그리고 그 문턱은 \emph{안쪽
수치만} 보고 놓을 수 있다(노트 498 의 $n_{\text{val}} \geq 150$ 은 유보
실패를 보고 놓았다). \textbf{그러나 우리는 지금 유보 결과를 알고 있으므로
이것을 검증이라 부르지 않는다. 다음 사이클에 사전 등록한다.}

\section{그래서 무엇이 문제였나}

고르개 셋을 나란히 놓으면 분명하다.

\begin{center}
\begin{tabular}{lrr}
\hline
고르개 & 뽑힌 짝 & 판 차 \\
\hline
도메인 점수(노트 498 $1$ 차) & $9$ & $-0.0085$ \\
안쪽 판(이 노트) & $5$ & $-0.0078$ \\
안쪽 판 $+$ 문턱 & $\mathbf{1}$ & $\mathbf{+0.0046}$ \\
\hline
\end{tabular}
\end{center}

\textbf{고르는 방법을 바꿔도 거의 안 바뀐다}($-0.0085 \to -0.0078$).
\textbf{바뀌는 것은 몇 개를 뽑느냐다.}

\begin{quote}
\textbf{조항} --- \textbf{안쪽 시기로 고른 것은 짝 수가 늘수록 나빠진다.}
안쪽 검증은 \emph{제일 센 짝 하나}는 짚지만 그 다음부터는 못 짚는다.
그러니 안쪽 고르개는 \textbf{순위}로만 쓰고 \textbf{개수는 따로
정한다}(하나부터 시작해 하나씩 늘리며 관문을 통과하는 데까지).
\end{quote}

\section*{남은 것}

\begin{itemize}
\item \textbf{안쪽 판 문턱 $0.003$ 을 사전 등록해 새 후보로 시험} ---
남은 축(target\_breadth · venue\_prominence · goods\_scale)은 아직 후보에
안 넣었다.
\item \textbf{왜 모바일 셋이 안 넘어가나} --- 노트 494 의 덮음 이동
인구조사에서 모바일이 네 짝으로 두 번째로 많다. 안쪽 시기와 유보 시기
사이에 모바일이 특히 많이 움직인다.
\item \textbf{안쪽 검증 시기를 옮겨 보기} --- $2023\!\sim\!2025$ 말고
$2022\!\sim\!2024$ 로 하면 다른 표가 나오나. 나오면 ``시기 탓''이 확정된다.
\end{itemize}

\end{document}
